\section{Marriage, Women, and the Distribution of Authority}

\subsection{Permission becomes failure}

Tertullian's thought about marriage cannot be represented by the radiant final
page of \emph{To His Wife} alone. That work's Christian couple share prayer,
fasting, psalm, persecution, and Eucharist in a genuinely reciprocal spiritual
life. Its purpose, however, is to prevent his widow from remarriage and at least
from marriage outside the Church. The household ideal is framed by his authority
to direct a wife's future after his death.

The later sequence makes continence progressively normative. In \emph{On
Exhortation to Chastity} 7--8 Tertullian knows Paul's permission of marriage but
argues that permission answers weakness and is not the highest good.
\emph{On Monogamy} treats one marriage as the Christian limit, even after a
spouse's death, under the Paraclete's final discipline. The argument elevates
fidelity and refuses to make bodies disposable. It also converts a charism of
continence into a rule capable of burdening widows and widowers whose second
marriages the wider Church permitted.

The change should not be psychologized from an undocumented household crisis.
Nothing records that his wife died, remarried, opposed him, or joined the New
Prophecy. The documents establish an addressed wife and an increasingly severe
sexual ethic. The domestic story between them is absent.

\subsection{Eve's daughters and the discipline of appearance}

\emph{On the Apparel of Women} 1.1 addresses Christian women collectively as
Eve's inheritors, the ``devil's gateway'' through whom the man was overcome and
the image of God destroyed. The rhetoric is not a modern mistranslation or a
harmless statement about Eve alone. Tertullian uses inherited blame to demand
mourning, modest dress, and renunciation of ornament.

The treatise also attacks an economy of display: precious stones, cosmetics,
hair, dye, and luxurious cloth carry labor, wealth, status competition, sexual
presentation, and myths of fallen angels. A social critique of conspicuous
consumption is present. It does not cancel the gendered humiliation by which
the critique is enforced. Tertullian can defend women's full bodily salvation
and accept women's prophetic speech while still making women bear a distinct
burden for temptation.

\emph{On the Veiling of Virgins} contests a Carthaginian custom by which
unmarried women remained unveiled in church. Tertullian argues from creation,
Paul, sexual maturity, other churches, nature, and the Paraclete that virgins as
women must be veiled. His opening slogan opposes truth to custom, but his own
case also appeals to custom when it supports the discipline. The work makes the
body publicly legible by marital and sexual status even as it claims the veil
protects inward holiness from display.

\subsection{Speaking prophets, barred ministers}

Chapter 9 of \emph{On the Veiling of Virgins} denies women the right to teach,
dispute, baptize, offer, or claim the function of a male office, much less the
priestly office. \emph{On Baptism} 17 similarly bars women from baptism despite
allowing laymen to act in necessity. These are direct restrictions, not merely
later church orders attributed to him.

At the same time, \emph{On the Soul} 9 grants evidentiary weight to a female
visionary who speaks after worship, and the New Prophecy's named founding
prophets include Prisca and Maximilla. Tertullian did not resolve the tension by
granting women equal institutional authority. Prophecy and office occupied
different channels: a woman might receive and report revelation while being
barred from teaching in church and from sacramental or presiding functions.
That distinction gives women real religious agency and keeps it under male
interpretation.

Modern narratives sometimes recruit him either as champion of charismatic
women or as an uncomplicated misogynist. Each isolates evidence. His works
permit women's inspired speech, construct the Christian marriage as a shared
discipline, condemn ornament economies, burden all baptized bodies with sexual
restraint, blame women through Eve, and exclude them from ministries he marks
male. The structure is complex without being egalitarian.

\subsection{Was the rigorist a priest?}

Jerome calls Tertullian a presbyter of the Church and then explains his later
break. No extant work says he was ordained. Two passages sound the other way.
In \emph{On Exhortation to Chastity} 7.3 he includes himself among the laity
while arguing that lay people too share priestly responsibility; in \emph{On
Monogamy} 12.2 he again opposes clergy and laity in language that places the
speaker with the latter. These are rhetorical and theological
self-locations, not civil certificates, but they make silent acceptance of
Jerome's late title unsafe.

Eusebius, who praises Tertullian and transmits the \emph{Apologeticum}, does not
call him presbyter. Cyprian's heavy literary debt does not preserve an office.
The responsible conclusion is not that ordination is impossible; it is that
priesthood rests on Jerome against suggestive authorial evidence and remains
doubtful.

The uncertainty matters because Tertullian writes imperiously without a secure
office. His authority comes from Scripture, rule of faith, rhetorical mastery,
ascetical credibility, and eventually prophecy. He can defend episcopal order
against lay presumption in \emph{On Baptism} and deny episcopal power in
\emph{On Modesty}. The movement is not necessarily a cleric surrendering his
office. It is an intellectual moving the criterion of true authority toward the
discipline he judges spiritual.
